\documentclass[12pt,a4paper]{article}

% Packages
\usepackage[a4paper, margin=1in]{geometry}
\usepackage{setspace}
\usepackage{times}
\usepackage{parskip}
\usepackage{booktabs}
\usepackage{array}
\usepackage{tabularx}
\usepackage{titlesec}
\usepackage{ragged2e}
\usepackage{microtype}
\usepackage[hidelinks]{hyperref}

% 1.5 line spacing
\onehalfspacing

% Section formatting - bold, 12pt, same as body
\titleformat{\section}{\normalfont\bfseries}{\thesection.}{0.5em}{}
\titleformat{\subsection}{\normalfont\bfseries}{\thesubsection}{0.5em}{}
\titlespacing*{\section}{0pt}{14pt}{6pt}
\titlespacing*{\subsection}{0pt}{10pt}{4pt}

% No paragraph indent, small space between paragraphs
\setlength{\parindent}{0pt}
\setlength{\parskip}{6pt}

\begin{document}

% ── COVER PAGE ───────────────────────────────────────────────────────────────
\begin{titlepage}
  \centering
  \vspace*{1.5cm}

  {\large\bfseries KYAMBOGO UNIVERSITY}\\[0.4cm]
  {\normalsize\bfseries FACULTY OF ENGINEERING}\\[0.4cm]
  {\normalsize\bfseries BACHELOR OF ENGINEERING IN MECHATRONICS AND BIOMEDICAL ENGINEERING}\\[0.4cm]
  {\normalsize\bfseries UNIVERSITY ELECTION MANAGEMENT SYSTEM}\\[0.2cm]
  {\normalsize\bfseries REPORT.}\\[1.2cm]

  \begin{tabular}{|>{\centering\arraybackslash\bfseries}p{6cm}|>{\centering\arraybackslash\bfseries}p{6cm}|}
    \hline
    NAME & REGISTRATION NUMBER \\
    \hline
    MATOVU ABDRHAMAN & 25/U/BIE/01392/PE \\
    \hline
    KWAGALA SOLOMON & 25/U/BIO/026/GV \\
    \hline
    MUGABI JOSHUA & 25/U/BIE/01394/PE \\
    \hline
    ALAN WANI VUCIRI & 25/U/BIO/024/GV \\
    \hline
    NAMUGANZA MARIA GORRET & 25/U/BIO/027/GV \\
    \hline
  \end{tabular}

\end{titlepage}

% ── 1. ABSTRACT ──────────────────────────────────────────────────────────────
\newpage
\section{ABSTRACT}

A University Election Management System was designed and implemented as a
console-based application in the C programming language to automate voter
registration, candidate management, vote casting, and result tabulation. The system
successfully prevented double voting using Boolean flags embedded in each voter
record and provided instant results upon completion of voting. All eight planned
features were implemented and verified within the defined constraints of 100 voters,
50 candidates, and 10 positions, demonstrating the effective application of C
programming concepts to a real-world administrative problem.

% ── 2. INTRODUCTION ──────────────────────────────────────────────────────────
\newpage
\section{INTRODUCTION}

Elections are a cornerstone of democratic governance in any institution. In a
university setting, student elections are critical to ensuring fair representation in
student bodies, guild councils, and academic committees. However, traditional manual
election processes are often slow, error-prone, and increasingly difficult to manage as
student populations grow. Issues such as voter fraud, double voting, and delayed
result announcements undermine the integrity and efficiency of the electoral process.

The University Election Management System was developed to address these challenges
by automating and streamlining the core processes involved in conducting a university
election. The system was built as a console-based application written in the C
programming language, chosen for its efficiency, direct memory access capabilities,
and suitability for demonstrating fundamental programming concepts including
struct-based data modelling, pointer manipulation, and array management.

The goal of this project was to design and implement a fully functional, menu-driven
election management system capable of handling position creation, voter and candidate
registration, record searching and updating, vote casting with fraud prevention, and
results display.

% ── 3. PROCEDURES AND OBSERVATIONS ──────────────────────────────────────────
\newpage
\section{PROCEDURES AND OBSERVATIONS}

\subsection{Program Design Structure}

The system was designed using a modular multi-file architecture to separate concerns
and improve maintainability. Six source files were created: \texttt{main.c} served as the
entry point containing the main menu loop; \texttt{electLibrary.h} declared all function
prototypes and system constants such as \texttt{MAX\_VOTERS} and \texttt{MAX\_CANDIDATES};
\texttt{electLibrary.c} contained the implementation of all core election functions;
\texttt{struct.h} defined the Voter and Candidate data structures; \texttt{arrays.h} declared
the global static arrays; and \texttt{readLine.h} together with \texttt{readLine.c} provided a
utility for safe string input handling. The application ran as a do-while loop in
\texttt{main.c}, continuously presenting the menu until the user selected the exit option,
with each menu choice dispatched to its corresponding function via a switch statement.
Two primary data structures were defined in \texttt{struct.h}. The Voter structure contained
an auto-assigned integer identifier, a fifty-character array for the voter's full name,
and a Boolean flag named \texttt{hasVoted} to track voting status. The Candidate structure
similarly contained an auto-assigned identifier, a name field, a position field, a
\texttt{hasVoted} flag, and an integer to accumulate the vote count.

\subsection{System Features}

The system was implemented with eight core features. Named election positions such as
President and Secretary could be created, with a maximum of ten positions supported.
New voters were enrolled with unique auto-generated identifiers, supporting up to one
hundred voters. Candidates were registered with a name and an assigned position, with
a maximum of fifty candidates. Voter and candidate records could be located by their
numeric identifier. A voter's name could be modified through an update function that
prompted for confirmation before saving changes. Voters could be removed from the
system with remaining records shifted to fill the gap. A registered voter could cast
exactly one vote for a candidate of their choice, with the system preventing any
subsequent votes from that voter. All candidates were displayed alongside their
positions and accumulated vote counts upon request.

\subsection{Important Functions}

The \texttt{registerVoter()} function registered new voters by checking capacity limits,
auto-assigning an identifier, setting \texttt{hasVoted} to false, and reading the voter name
through \texttt{readLine()} to prevent buffer overflow. The \texttt{castVote()} function located
the voter using \texttt{findVoter()} and, if the voter had not yet voted, accepted a candidate
identifier, located the candidate using \texttt{findCandidate()}, incremented the candidate's
vote count, and permanently set the voter's \texttt{hasVoted} flag to true. The
\texttt{deleteVoter()} function used pointer arithmetic to compute the array index of the
target voter and shifted all subsequent records one position to the left before
decrementing the voter count. The \texttt{findVoter()} and \texttt{findCandidate()} helper
functions performed linear searches and returned a pointer to the matching struct or
NULL if not found, with all calling functions verified to handle the NULL case. The
\texttt{isValidPosition()} function used \texttt{strcmp()} to validate uniqueness before adding a
position. The \texttt{readLine()} function flushed the input buffer, read a full line using
\texttt{fgets()}, and stripped the trailing newline using \texttt{strcspn()} to eliminate the
common C input corruption issue.

\subsection{Design Decisions and Justifications}

Fixed-size static arrays were chosen over dynamically allocated structures to keep the
implementation simple and avoid memory leak risks associated with \texttt{malloc} and
\texttt{free}. For a university election context with defined upper limits, this approach was
appropriate and sufficient. The \texttt{hasVoted} boolean embedded directly in each Voter
struct provided an O(1) check at vote time, more efficient than maintaining a separate
registry of voted voters. Voter and candidate identifiers were assigned as count plus
one at registration, eliminating user errors and guaranteeing uniqueness within a
session. The codebase was split across six files to improve readability, enable parallel
development, and simplify individual function testing. A dedicated \texttt{readLine()} utility
was implemented in place of standard \texttt{scanf} or \texttt{gets} to guarantee safe and
predictable string input, representing a deliberate defensive design choice against
common C input vulnerabilities.

% ── 4. DISCUSSION ────────────────────────────────────────────────────────────
\newpage
\section{DISCUSSION}

The University Election Management System achieved all of its stated goals,
demonstrating that fundamental C programming constructs can effectively automate a
real-world administrative process. The significance of these findings is that a
structured, menu-driven application with well-defined data structures and disciplined
input handling is sufficient to eliminate the most critical weaknesses of manual
election management, namely voter fraud, tallying errors, and lack of transparency in
results.

In comparison with theoretical expectations, the system performed exactly as designed.
The Boolean flag mechanism for double-vote prevention functioned with the expected
O(1) efficiency, and all linear search functions performed within acceptable bounds
given the fixed array sizes. No deviation from the intended logical behaviour was
observed during testing, meaning the implemented system aligned fully with its
specification.

The absence of position validation in \texttt{registerCandidate()} represented a source of
systematic error, as a candidate could be assigned to a non-existent position without
any system warning, potentially producing inconsistent election data. The unsorted
output of \texttt{displayResults()} reduced result clarity but did not affect correctness. The
primary systematic error encountered during development was the input buffer conflict
from mixed use of \texttt{scanf} and \texttt{fgets}, which consistently produced blank name entries
and was fully resolved through the custom \texttt{readLine()} function. No random errors were
observed, as the system operated deterministically within its defined constraints.

The procedure could be improved by implementing position validation in
\texttt{registerCandidate()} using \texttt{strcmp()} against the registered positions array, adding
a sort by vote count before displaying results, and introducing file input and output to
enable persistent storage across sessions. In terms of practical applications, this
system could serve any institution requiring a lightweight offline election management
tool. With the addition of file I/O, sorted results, position validation, and a graphical
user interface, it could function as a fully featured standalone election application for
small to medium-sized institutions.

\subsection{Challenges Encountered}

Mixing \texttt{scanf} for integers with \texttt{fgets} for strings caused the leftover newline from
\texttt{scanf} to be immediately consumed by \texttt{fgets}, resulting in blank name entries. This
was resolved by the custom \texttt{readLine()} function, which always flushes the input buffer
before reading. Correctly removing a voter from the middle of the array using pointer
arithmetic required careful understanding of C struct pointer behaviour, specifically
computing the index as the difference between the voter pointer and the base of the
voters array before shifting. The absence of persistent storage meant that all data was
lost when the program exited, with file input and output identified as a necessary
future improvement. The \texttt{displayResults()} function listed candidates in registration
order rather than ranked by votes, and adding a sort before display was identified as a
target for future development.

\subsection{Problems Solved}

The system resolved the key problems of manual university election management.
Double voting and voter fraud were prevented through the \texttt{hasVoted} Boolean flag,
providing a reliable and efficient programmatic check. The error-prone manual tallying
of votes was replaced by automated vote counting within \texttt{castVote()}. Structural
management of candidate registrations and position assignments was achieved through
the defined data structures and menu-driven interface. Instant and transparent result
display was made possible by the \texttt{displayResults()} function. The inability to update
or remove invalid voter records in a manual system was addressed by the implemented
update and delete functions, both of which operated correctly and safely within the
array-based data model.

% ── 5. CONCLUSION ────────────────────────────────────────────────────────────
\newpage
\section{CONCLUSION}

The University Election Management System was successfully designed and implemented
as a modular, console-based C application that met all stated project goals by
automating voter and candidate registration, enforcing double-vote prevention, and
providing instant result display. The project demonstrated the practical application of
struct-based data modelling, pointer manipulation, array management, and defensive
input handling in C, and established a solid foundation for future extension with
persistent storage, sorted results, and a graphical interface.

% ── 6. REFERENCES ────────────────────────────────────────────────────────────
\newpage
\section{REFERENCES}

\begin{enumerate}
  \item Kernighan, B. W., \& Ritchie, D. M. (1988). \textit{The C Programming Language} (2nd ed.). Prentice Hall.
  \item Deitel, P., \& Deitel, H. (2016). \textit{C How to Program} (8th ed.). Pearson Education.
  \item King, K. N. (2008). \textit{C Programming: A Modern Approach} (2nd ed.). W. W. Norton \& Company.
\end{enumerate}

\end{document}
